\documentclass[11pt,a4paper]{article}
\usepackage[french]{babel}
\usepackage[utf8]{inputenc}
\usepackage[T1]{fontenc}
\usepackage{amsmath}
\usepackage{amsthm}
\usepackage{amsfonts}
\usepackage{amssymb}
\usepackage[left=2cm, right=2cm, top=2cm, bottom=2cm]{geometry}
\usepackage{hyperref}
\usepackage{mdframed}
\usepackage{xcolor}

\author{Ulrich TCHISSAMBOU}
\title{\textbf{Démonstrations}\\[4pt]
       \large Mathématiques --- Terminale Spécialité}
\date{}

\newcommand{\vc}[1]{\overrightarrow{#1}}

\newtheoremstyle{enoncestyle}
  {8pt}{8pt}{\normalfont}{0pt}{\bfseries\color{blue!70!black}}{.}{.5em}{}
\theoremstyle{enoncestyle}
\newtheorem{enonce}{Énoncé}[section]

\renewenvironment{proof}[1][\proofname]{%
  \par\medskip
  \noindent\textbf{\color{black!70}#1.}\quad\itshape
}{%
  \hfill$\blacksquare$\par\medskip
}

\setcounter{tocdepth}{2}

\begin{document}
\maketitle
\tableofcontents
\newpage

% ============================================================
\section{Chapitre 2 — Limites de suite}
% ============================================================

\subsection{Unicité de la limite d'une suite}

\begin{mdframed}[linecolor=blue!40,backgroundcolor=blue!4,linewidth=1pt]
\begin{enonce}
Si une suite $(u_n)$ a une limite réelle, alors cette limite est unique.
\end{enonce}
\end{mdframed}

\begin{proof}[Démonstration]
Supposons que $(u_n)$ ait deux limites $\ell_1$ et $\ell_2$ avec $\ell_1 \neq \ell_2$.
Posons :
\[
    \varepsilon = \frac{|\ell_1 - \ell_2|}{3} > 0.
\]
Par définition de la limite :
\begin{itemize}
    \item il existe $N_1$ tel que pour tout $n \geq N_1$, $|u_n - \ell_1| < \varepsilon$ ;
    \item il existe $N_2$ tel que pour tout $n \geq N_2$, $|u_n - \ell_2| < \varepsilon$.
\end{itemize}
Pour $n \geq \max(N_1, N_2)$, par inégalité triangulaire :
\[
    |\ell_1 - \ell_2| \leq |u_n - \ell_1| + |u_n - \ell_2|
    < 2\varepsilon = \frac{2}{3}|\ell_1 - \ell_2|.
\]
On obtient $|\ell_1 - \ell_2| < \dfrac{2}{3}|\ell_1 - \ell_2|$, ce qui est impossible
si $\ell_1 \neq \ell_2$. Donc $\ell_1 = \ell_2$.
\end{proof}

% ============================================================
\section{Chapitre 3 — Vecteurs, droites et plans de l'espace}
% ============================================================

\subsection{Parallélisme de l'intersection de deux plans}

\begin{mdframed}[linecolor=blue!40,backgroundcolor=blue!4,linewidth=1pt]
\begin{enonce}
Si deux plans $\mathcal{P}$ et $\mathcal{P}'$ contiennent respectivement deux droites
$(d)$ et $(d')$ parallèles entre elles, alors leur intersection
$(d'') = \mathcal{P} \cap \mathcal{P}'$ est parallèle à $(d)$ et $(d')$.
\end{enonce}
\end{mdframed}

\begin{proof}[Démonstration]
Soit $\vc{v}$ un vecteur directeur commun à $(d)$ et $(d')$ (elles sont parallèles, donc
colinéaires).

Soit $\vc{w}$ un vecteur directeur de la droite d'intersection
$(d'') = \mathcal{P} \cap \mathcal{P}'$.

Supposons par l'absurde que $\vc{w}$ et $\vc{v}$ ne soient pas colinéaires. Alors
$(\vc{v}, \vc{w})$ serait une base commune à $\mathcal{P}$ (qui contient $(d)$ de vecteur
$\vc{v}$ et $(d'')$ de vecteur $\vc{w}$) et à $\mathcal{P}'$ (qui contient $(d')$ de
vecteur $\vc{v}$ et $(d'')$ de vecteur $\vc{w}$). Cela impliquerait
$\mathcal{P} = \mathcal{P}'$, contrairement à l'hypothèse qu'ils sont sécants.

Donc $\vc{w}$ et $\vc{v}$ sont colinéaires : $(d'')$ est parallèle à $(d)$ et $(d')$.
\end{proof}

% ============================================================
\section{Chapitre 8 — Loi binomiale}
% ============================================================

\subsection{Nombre de parties d'un ensemble fini}

\begin{mdframed}[linecolor=blue!40,backgroundcolor=blue!4,linewidth=1pt]
\begin{enonce}
Le nombre de parties d'un ensemble fini à $n$ éléments est $2^n$.
\end{enonce}
\end{mdframed}

\begin{proof}[Démonstration]
Soit $E = \{x_1, x_2, \ldots, x_n\}$. À toute partie $A \subseteq E$, on associe le mot
binaire $(b_1, b_2, \ldots, b_n)$ défini par :
\[
    b_i = \begin{cases} 1 & \text{si } x_i \in A, \\ 0 & \text{sinon.}\end{cases}
\]
Cette correspondance est une bijection entre les parties de $E$ et les mots binaires de
longueur $n$. Chaque $b_i$ pouvant prendre $2$ valeurs indépendamment, il y a :
\[
    \underbrace{2 \times 2 \times \cdots \times 2}_{n \text{ facteurs}} = 2^n
\]
mots distincts, donc $2^n$ parties.
\end{proof}

\subsection{Formule de la loi binomiale $P(X=k)$}

\begin{mdframed}[linecolor=blue!40,backgroundcolor=blue!4,linewidth=1pt]
\begin{enonce}
Soit $X \sim \mathcal{B}(n;p)$. Pour tout entier $k \in \{0, 1, \ldots, n\}$ :
\[
    P(X = k) = \binom{n}{k}\,p^k\,(1-p)^{n-k}.
\]
\end{enonce}
\end{mdframed}

\begin{proof}[Démonstration]
On répète $n$ épreuves de Bernoulli indépendantes, chaque succès ayant probabilité $p$.

\textbf{Probabilité d'un chemin fixé.}
Un chemin réalisant $k$ succès dans des positions déterminées a pour probabilité :
\[
    \underbrace{p \cdots p}_{k}\; \times \underbrace{(1-p)\cdots(1-p)}_{n-k}
    = p^k(1-p)^{n-k}.
\]

\textbf{Nombre de chemins favorables.}
Choisir les positions des $k$ succès parmi $n$ épreuves revient à choisir une partie à $k$
éléments de $\{1, \ldots, n\}$, soit $\dbinom{n}{k}$ chemins distincts.

\textbf{Conclusion.}
Ces chemins sont deux à deux incompatibles, donc par la règle d'additivité :
\[
    P(X = k) = \binom{n}{k}\,p^k\,(1-p)^{n-k}.
\]
\end{proof}

% ============================================================
\section{Chapitre 10 — Primitives et équations différentielles}
% ============================================================

\subsection{Solutions de $y' = ay$}

\begin{mdframed}[linecolor=blue!40,backgroundcolor=blue!4,linewidth=1pt]
\begin{enonce}
Les solutions de l'équation différentielle $y' = ay$ (avec $a \in \mathbf{R}^*$)
sont exactement les fonctions $x \mapsto k\,e^{ax}$, $k \in \mathbf{R}$.
\end{enonce}
\end{mdframed}

\begin{proof}[Démonstration]
\textbf{Ces fonctions sont bien solutions.}
Posons $f_k(x) = ke^{ax}$. Alors :
\[
    f_k'(x) = k\,a\,e^{ax} = a\,f_k(x),
\]
donc $f_k$ vérifie $y' = ay$.

\textbf{Ce sont les seules solutions.}
Soit $y$ une solution quelconque. Posons $g(x) = y(x)\,e^{-ax}$. Sa dérivée est :
\begin{align*}
    g'(x) &= y'(x)\,e^{-ax} + y(x)\cdot(-a)\,e^{-ax} \\
          &= a\,y(x)\,e^{-ax} - a\,y(x)\,e^{-ax} = 0.
\end{align*}
$g$ est donc constante : $g(x) = k$ pour un certain $k \in \mathbf{R}$, ce qui donne
$y(x) = k\,e^{ax}$.
\end{proof}

\subsection{Solutions de $y' = ay + b$}

\begin{mdframed}[linecolor=blue!40,backgroundcolor=blue!4,linewidth=1pt]
\begin{enonce}
Les solutions de l'équation différentielle $y' = ay + b$ ($a \neq 0$, $b \neq 0$)
sont les fonctions $x \mapsto k\,e^{ax} - \dfrac{b}{a}$, $k \in \mathbf{R}$.
\end{enonce}
\end{mdframed}

\begin{proof}[Démonstration]
\textbf{Étape 1 : solution particulière constante.}
Cherchons $f : x \mapsto c$ solution. Alors $f'= 0$, et l'équation donne $0 = ac + b$,
soit $c = -\dfrac{b}{a}$. La fonction $f : x \mapsto -\dfrac{b}{a}$ est une solution
particulière.

\textbf{Étape 2 : résoudre l'équation homogène.}
Soit $g$ une solution quelconque de $y' = ay + b$. Posons $h = g - f$. Alors :
\[
    h'(x) = g'(x) - f'(x)
    = \bigl(a\,g(x) + b\bigr) - \bigl(a\,f(x) + b\bigr)
    = a\,\bigl(g(x) - f(x)\bigr) = a\,h(x).
\]
Donc $h$ est solution de $y' = ay$, d'où $h(x) = k\,e^{ax}$ (démonstration précédente).

\textbf{Conclusion.}
\[
    g(x) = h(x) + f(x) = k\,e^{ax} - \frac{b}{a}.
\]
\end{proof}

% ============================================================
\section{Chapitre 11 — Combinatoire et dénombrement}
% ============================================================

\subsection{Nombre de parties d'un ensemble fini}

La démonstration est identique à celle du chapitre~8, section~3.1 :
à chaque partie de $E$ correspond un unique mot binaire de longueur $n$, et il y a
$2^n$ tels mots. \hfill$\blacksquare$

\end{document}
